\documentclass{amsart}
\usepackage{amsmath,amssymb,amsthm,hyperref}
\theoremstyle{plain}
\newtheorem{theorem}{Theorem}
\newtheorem{proposition}{Proposition}
\newtheorem{lemma}{Lemma}
\theoremstyle{remark}
\newtheorem{remark}{Remark}
\newtheorem*{statusnote}{Status}
\DeclareMathOperator{\rad}{rad}
\begin{document}

\title{On the equation $x(x+1)=y!$}
\maketitle

\begin{abstract}
We study the Diophantine equation $x(x+1)=y!$ in positive integers.
The only solutions presently known are $(x,y)=(1,2)$ and $(x,y)=(2,3)$,
and we verify computationally that there are no others with $y\le 2000$.
We prove rigorously every fact that the current state of the art allows:
the reformulation of the equation as ``$4\,y!+1$ is a perfect square'',
the coprime two-factor structure, a complete elimination of the small
cases $y\le 7$ by an explicit pronic-interval argument, sharp lower bounds
on the two factors, and the fact (Berend--Osgood) that the equation has
only finitely many solutions.  We also prove three precise \emph{negative}
metatheorems showing that the elementary obstructions which settle most
factorial Diophantine equations -- congruences, prime-valuation parity, and
size of fixed prime powers -- are provably powerless here.  We state honestly
that a \emph{complete} determination of the solution set is, at present, an
\textbf{open problem} of Brocard--Ramanujan type: the finiteness theorem is
ineffective and no effective bound on $y$ is known, so we cannot exclude a
further (necessarily enormous) solution. We make explicit exactly which
ingredient is missing.
\end{abstract}

\section{Statement and summary}

We consider
\begin{equation}\label{eq:main}
   x(x+1)=y!,\qquad x,y\in\mathbb{Z}_{\ge 1},
\end{equation}
where $y!=1\cdot 2\cdots y$.

\begin{statusnote}
The task of \emph{determining all} solutions of \eqref{eq:main} is, to the
best of current mathematical knowledge, \textbf{open}.  The only solutions
known are
\[
   (x,y)=(1,2)\quad\bigl(1\cdot 2=2!\bigr),\qquad
   (x,y)=(2,3)\quad\bigl(2\cdot 3=3!\bigr),
\]
and exhaustive computation (Section~\ref{sec:comput}) shows there is no
further solution with $y\le 2000$.  It is a theorem of Berend--Osgood
(Section~\ref{sec:finite}) that \eqref{eq:main} has only finitely many
solutions, but that theorem is \emph{ineffective}: it produces no explicit
bound on $y$, and therefore does \emph{not} entail that the list above is
complete.  We do not claim a complete proof of completeness; instead we
prove, with full rigour, every statement below, and we establish three
metatheorems (Section~\ref{sec:obstruction}) showing precisely why the
elementary techniques that suffice for many factorial Diophantine equations
are insufficient here. The missing ingredient is identified in
Remark~\ref{rem:missing}.
\end{statusnote}

The remainder of the paper records exactly what can be proved rigorously,
organised so that each section is self-contained.

\section{Reformulation and elementary structure}

\begin{proposition}\label{prop:square}
A positive integer $y$ admits a solution $x$ of \eqref{eq:main} if and only
if $4\,y!+1$ is a perfect square.  In that case the solution is unique and
given by $x=\tfrac{1}{2}\bigl(\sqrt{4\,y!+1}-1\bigr)$.
\end{proposition}

\begin{proof}
Equation \eqref{eq:main} is the quadratic $x^2+x-y!=0$ in $x$, whose
discriminant is $1+4\,y!$.  Its two real roots are
$\frac{-1\pm\sqrt{1+4\,y!}}{2}$, and only the root with the $+$ sign is
positive.  Hence a positive root lies in $\mathbb{Z}_{\ge1}$ iff
$\sqrt{1+4\,y!}\in\mathbb{Z}$, i.e.\ iff $1+4\,y!$ is a perfect square: in
that case $1+4\,y!$ is odd, so $\sqrt{1+4\,y!}$ is an odd integer,
$\sqrt{1+4\,y!}-1$ is even, and
$x=\frac{-1+\sqrt{1+4\,y!}}{2}\in\mathbb{Z}_{\ge 1}$ (positivity holds
because $y\ge1$ gives $4\,y!+1\ge5$, so $\sqrt{4\,y!+1}\ge\sqrt5>1$).
Conversely if $x$ solves \eqref{eq:main} then
$(2x+1)^2=4x^2+4x+1=4\,y!+1$ is a perfect square. Uniqueness of $x$ for each
$y$ is immediate since $x\mapsto x(x+1)$ is strictly increasing on
$\mathbb{Z}_{\ge 1}$.
\end{proof}

Thus \eqref{eq:main} is equivalent to the requirement that $4\,y!+1$ be a
square; with $m=2x+1$ it reads
\begin{equation}\label{eq:m}
   m^2=4\,y!+1,\qquad m\ \text{odd}.
\end{equation}
This places \eqref{eq:main} in the Brocard--Ramanujan family
``$c\cdot n!+1=\square$''; the classical Brocard problem is the case
$c=1$, $n!+1=\square$, which is itself famously open.

\begin{lemma}\label{lem:coprime}
If $x(x+1)=y!$ then $\gcd(x,x+1)=1$, so $x$ and $x+1$ are coprime divisors
of $y!$ whose product is $y!$.  Consequently every prime $p\le y$ satisfies
$p^{\,v_p(y!)}\mid x$ or $p^{\,v_p(y!)}\mid x+1$ (the entire $p$-part of $y!$
lies in a single factor), and one has the explicit description
\begin{equation}\label{eq:coprimedesc}
   x=\prod_{p\in S}p^{\,v_p(y!)},\qquad
   x+1=\prod_{p\le y,\;p\notin S}p^{\,v_p(y!)},
\end{equation}
for a uniquely determined set $S$ of primes $\le y$, where $v_p(\cdot)$
denotes the $p$-adic valuation.
\end{lemma}

\begin{proof}
Any common divisor of $x$ and $x+1$ divides their difference $1$, so
$\gcd(x,x+1)=1$.  Fix a prime $p\le y$ and let $a=v_p(y!)\ge1$.  Since
$x(x+1)=y!$, we have $v_p(x)+v_p(x+1)=a$; but $\gcd(x,x+1)=1$ forces
$\min\{v_p(x),v_p(x+1)\}=0$, so exactly one of $v_p(x),v_p(x+1)$ equals $a$
and the other is $0$.  Letting $S$ be the set of primes $p\le y$ with
$v_p(x)=a$ gives \eqref{eq:coprimedesc}: indeed for such $p$ the full power
$p^a$ divides $x$ and not $x+1$, and for $p\notin S$ the full power $p^a$
divides $x+1$ and not $x$. Since the only primes dividing $y!$ are those
$\le y$, the two displayed products are coprime with product
$\prod_{p\le y}p^{v_p(y!)}=y!$, and $S$ is determined by $x$.
\end{proof}

\section{Elimination of the small cases $y\le 7$}\label{sec:small}

We dispose of the small range by a completely explicit, hand-verifiable
argument that does not rely on any computation beyond multiplying small
integers.  The idea: a positive integer $N$ is \emph{pronic} (of the form
$a(a+1)$) iff it equals one of the strictly increasing pronic numbers
$0,2,6,12,20,30,\dots$; if $N$ lies strictly between two consecutive pronic
numbers then $N$ is not pronic, so \eqref{eq:main} has no solution with
$y!=N$.

\begin{proposition}\label{prop:small}
The only solutions of \eqref{eq:main} with $1\le y\le 7$ are $(x,y)=(1,2)$
and $(x,y)=(2,3)$.
\end{proposition}

\begin{proof}
The function $a\mapsto a(a+1)$ is strictly increasing on $\mathbb{Z}_{\ge0}$,
so if $a(a+1)<N<(a+1)(a+2)$ for some integer $a\ge0$, then no integer $b$
satisfies $b(b+1)=N$ (any such $b\ge0$ would give $a<b<a+1$, impossible),
and \eqref{eq:main} has no positive solution with $y!=N$.

For $y=1$: $1!=1$, and $x(x+1)\ge 1\cdot2=2$ for all $x\ge1$, so there is no
positive solution.

For $y=2$: $2!=2=1\cdot2$, giving $x=1$.

For $y=3$: $3!=6=2\cdot3$, giving $x=2$.

For $y=4$: $4!=24$, and $4\cdot5=20<24<30=5\cdot6$. No solution.

For $y=5$: $5!=120$, and $10\cdot11=110<120<132=11\cdot12$. No solution.

For $y=6$: $6!=720$, and $26\cdot27=702<720<756=27\cdot28$. No solution.

For $y=7$: $7!=5040$, and $70\cdot71=4970<5040<5112=71\cdot72$. No solution.

In each non-solution case $y!$ is strictly trapped between two consecutive
pronic numbers, hence is itself not pronic, so \eqref{eq:main} fails.  All
the displayed products are elementary integer multiplications.
\end{proof}

\section{Lower bounds on the factors}\label{sec:bounds}

The following bounds, proved unconditionally, sharpen the structural picture
and are used in Section~\ref{sec:obstruction} to explain why size arguments
cannot obstruct large solutions.

\begin{lemma}\label{lem:size}
If $x(x+1)=y!$ with $y\ge1$, then
\[
   \sqrt{y!}-1<x<\sqrt{y!},\qquad \sqrt{y!}<x+1<\sqrt{y!}+1 .
\]
In particular both factors $x,\ x+1$ exceed $\sqrt{y!}-1$.
\end{lemma}

\begin{proof}
From $x(x+1)=y!$ we get $x^2<x(x+1)=y!<(x+1)^2$, so $x<\sqrt{y!}<x+1$, i.e.\
$x+1>\sqrt{y!}$ and $x<\sqrt{y!}$. Combining $x<\sqrt{y!}$ with
$x+1>\sqrt{y!}$ (so $x>\sqrt{y!}-1$) and $x+1<\sqrt{y!}+1$ (since
$x<\sqrt{y!}$) gives the claim.
\end{proof}

\begin{lemma}\label{lem:factbig}
For every integer $y\ge4$ one has $y!>y^2$, hence $\sqrt{y!}>y$.
\end{lemma}

\begin{proof}
For $y\ge4$ the product $y!$ contains the three distinct factors $2$, $y-1$,
$y$ (distinct because $y-1\ge3>2$), so $y!\ge 2\,(y-1)\,y$. Since $y\ge3$
gives $2(y-1)>y$, we obtain $y!\ge 2(y-1)y>y\cdot y=y^2$. Taking square
roots gives $\sqrt{y!}>y$.
\end{proof}

\begin{lemma}\label{lem:bigprime}
Let $y\ge4$ and let $p$ be any prime with $y/2<p\le y$ (such a prime exists
by Bertrand's postulate). If $x(x+1)=y!$ then $v_p(y!)=1$, and $p$ divides
exactly one of $x,\,x+1$. Moreover $p\le y<\sqrt{y!}$, so $p$ alone is far
too small to account for the size $>\sqrt{y!}-1$ of the factor it divides.
\end{lemma}

\begin{proof}
By Legendre's formula $v_p(y!)=\sum_{j\ge1}\lfloor y/p^j\rfloor$. Since
$y/2<p\le y$ we have $1\le y/p<2$, so $\lfloor y/p\rfloor=1$, while
$p^2>(y/2)^2\ge y$ for $y\ge4$ gives $\lfloor y/p^j\rfloor=0$ for $j\ge2$;
hence $v_p(y!)=1$. By Lemma~\ref{lem:coprime} the full $p$-part $p^1=p$ lies
in exactly one factor. Finally $p\le y<\sqrt{y!}$ by Lemma~\ref{lem:factbig},
and by Lemma~\ref{lem:size} the factor of $x,x+1$ that $p$ divides exceeds
$\sqrt{y!}-1$, which (for $y\ge4$, where $\sqrt{y!}>y\ge p$) is larger than
$p-1$; thus $p$ cannot by itself fill out that factor.
\end{proof}

\section{Why the standard elementary obstructions fail}\label{sec:obstruction}

We record, as rigorous \emph{metatheorems}, that the three techniques which
most often resolve factorial equations are provably powerless here.  These
are not failed attempts but theorems about the futility of certain methods;
they delineate exactly why \eqref{eq:main} is hard.

\begin{proposition}[No congruence obstruction]\label{prop:nocong}
Let $N\ge2$ be any fixed modulus. For every $y\ge N$ we have
$4\,y!+1\equiv1\pmod N$, and $1$ is a quadratic residue modulo every modulus
(being $1\equiv 1^2$). Hence the necessary local condition for solubility of
\eqref{eq:m} -- that $4\,y!+1$ be a square modulo $N$ -- holds for every
fixed $N$ once $y\ge N$. In particular no single congruence condition can
exclude solutions of \eqref{eq:main} for all large $y$.
\end{proposition}

\begin{proof}
If $y\ge N$ then $N\mid y!$, because $N\le y$ implies $N$ is one of the
integer factors $1,2,\dots,y$ whose product is $y!$; hence
$4\,y!\equiv0\pmod N$ and $4\,y!+1\equiv1\pmod N$. As $1\equiv1^2\pmod N$,
the congruence $m^2\equiv4\,y!+1\pmod N$ is solvable (take $m\equiv1$). Thus
the necessary local condition for \eqref{eq:m} holds modulo every $N$ once
$y\ge N$. Since any congruence-based exclusion would have to fail this local
condition at some modulus, no such exclusion is available for $y\ge N$.
\end{proof}

\begin{proposition}[No valuation-parity obstruction]\label{prop:noparity}
The parity of the exponents $v_p(y!)$ imposes \emph{no} constraint that can
exclude a solution. Concretely: Lemma~\ref{lem:coprime} forces the whole
$p$-part $p^{v_p(y!)}$ into a single factor $x$ or $x+1$, but it never forces
either factor to be a perfect square, and a genuine solution exists with
$v_p(y!)$ odd for some $p$.
\end{proposition}

\begin{proof}
A common heuristic for factorial equations is: if $v_p(y!)$ is odd then ``a
factor is forced to be a non-square,'' producing a contradiction with some
squareness requirement. No such requirement exists here. Equation
\eqref{eq:main} only requires $x(x+1)=y!$; neither $x$ nor $x+1$ need be a
square, and \eqref{eq:m} requires $4\,y!+1$ (not $y!$ or its factors) to be
a square. By Lemma~\ref{lem:coprime}, an odd $v_p(y!)$ merely places the odd
prime power $p^{v_p(y!)}$ inside one of $x,x+1$; the complementary primes can
freely supply whatever further factors are needed, with no squareness forced
on either side. Indeed, in the genuine solution $2\cdot3=3!$ we have
$v_2(3!)=1$ and $v_3(3!)=1$, both \emph{odd}, yet $(x,y)=(2,3)$ is a
solution. Hence no argument based solely on the parity of the exponents
$v_p(y!)$ can preclude solutions of \eqref{eq:main}.
\end{proof}

\begin{proposition}[No fixed-prime-power size obstruction]\label{prop:nosize}
Fix any integer $k\ge1$ and any primes $p_1,\dots,p_k\le y$, and let
$Q=\prod_{i=1}^{k}p_i^{\,v_{p_i}(y!)}$ be the partial product that, by
Lemma~\ref{lem:coprime}, divides a single factor $A\in\{x,x+1\}$ whenever
$\{p_1,\dots,p_k\}$ lies entirely in $S$ or entirely in its complement.
The forced inequality $A\ge Q$ cannot contradict the size bound
$A\le\sqrt{y!}+1$ of Lemma~\ref{lem:size}, because for every fixed $k$ and
every choice of primes,
\[
   \frac{Q}{\sqrt{y!}}\longrightarrow 0\qquad(y\to\infty).
\]
\end{proposition}

\begin{proof}
We first bound each prime power dividing $y!$ by the largest one, $2^{v_2(y!)}$.

\emph{Step 1: $v_p(y!)\le y/(p-1)$ for every prime $p$.} By Legendre's
formula and the geometric series,
\[
   v_p(y!)=\sum_{j\ge1}\Big\lfloor \tfrac{y}{p^j}\Big\rfloor
   \le\sum_{j\ge1}\tfrac{y}{p^j}=\frac{y}{p-1}.
\]

\emph{Step 2: $2^{v_2(y!)}$ is the maximal prime power, and $\le 2^{y}$.}
From Step 1 with $p=2$, $v_2(y!)\le y$, so $2^{v_2(y!)}\le 2^{y}$. For an
odd prime $p\ge3$, Step 1 gives $p^{v_p(y!)}\le p^{\,y/(p-1)}$. The function
$t\mapsto t^{1/(t-1)}$ is decreasing for $t\ge2$, so for $p\ge3$ we have
$p^{1/(p-1)}\le 3^{1/2}<2$, whence
$p^{\,y/(p-1)}=\bigl(p^{1/(p-1)}\bigr)^{y}<2^{y}$. Therefore every prime
power $p^{v_p(y!)}$ with $p\le y$ satisfies $p^{v_p(y!)}\le 2^{y}$.

\emph{Step 3: bound on $Q$.} Using Step 2 for each of the $k$ factors,
\[
   Q=\prod_{i=1}^{k}p_i^{\,v_{p_i}(y!)}\le\bigl(2^{y}\bigr)^{k}=2^{ky},
   \qquad\text{so}\qquad \log Q\le k\,y\log 2 .
\]

\emph{Step 4: comparison with $\sqrt{y!}$.} By Stirling's formula,
\[
   \log\sqrt{y!}=\tfrac12\log(y!)
   =\tfrac12\bigl(y\log y-y+O(\log y)\bigr)\sim\tfrac12\,y\log y .
\]
Hence
\[
   \log\!\Big(\frac{Q}{\sqrt{y!}}\Big)\le k\,y\log 2-\tfrac12\,y\log y
   = y\Big(k\log 2-\tfrac12\log y\Big)\longrightarrow-\infty
   \qquad(y\to\infty),
\]
because $\tfrac12\log y>k\log 2$ once $y>2^{2k}$. Therefore
$Q/\sqrt{y!}\to0$, and for all large $y$ the constraint $A\ge Q$ is
consistent with $A\le\sqrt{y!}+1$. Thus no obstruction can arise from
forcing finitely many prime powers into one factor and bounding by size.
\end{proof}

Propositions~\ref{prop:nocong}--\ref{prop:nosize} together explain, in a
precise quantitative sense, why \eqref{eq:main} resists every elementary
attack: there is no congruence barrier, no parity barrier, and no
fixed-prime-power size barrier. The factors $x,x+1\approx\sqrt{y!}$ are far
larger than any bounded combination of primes $\le y$
(Lemma~\ref{lem:bigprime}, Proposition~\ref{prop:nosize}), leaving the prime
powers ample freedom to distribute across the two coprime factors.

\section{Finiteness of the solution set}\label{sec:finite}

\begin{theorem}[Berend--Osgood]\label{thm:finite}
Let $P\in\mathbb{Z}[t]$ have degree $\ge 2$.  Then the equation $P(x)=y!$
has only finitely many solutions $(x,y)\in\mathbb{Z}_{\ge 1}^2$.  In
particular, $x(x+1)=y!$ has only finitely many solutions.
\end{theorem}

This is the main result of D.~Berend and C.~Osgood, \emph{On the equation
$P(x)=n!$ and a question of Erd\H{o}s}, J.\ Number Theory \textbf{42}
(1992), 189--193.  Applying it with $P(t)=t(t+1)=t^2+t$ (degree $2\ge2$)
gives the finiteness of the solution set of \eqref{eq:main}. We recall its
meaning and limitation rather than reproducing its proof.

\begin{remark}[The theorem is ineffective]\label{rem:ineff}
The Berend--Osgood argument shows the solution set is finite but does
\emph{not} produce an explicit bound on $y$. It therefore does \emph{not}
imply that $(1,2)$ and $(2,3)$ exhaust the solutions: it leaves open the
logical possibility of a further (necessarily enormous) solution. This is
precisely why the determination problem remains open even though finiteness
is a theorem. No effective upper bound for $y$ in \eqref{eq:main} is known.
\end{remark}

\section{Computational verification}\label{sec:comput}

\begin{proposition}\label{prop:comp}
For $1\le y\le 2000$ the only $y$ for which $4\,y!+1$ is a perfect square
are $y=2$ and $y=3$.  Hence the only solutions of \eqref{eq:main} with
$y\le 2000$ are $(1,2)$ and $(2,3)$.
\end{proposition}

\begin{proof}[Verification]
By Proposition~\ref{prop:square} it suffices to test, for each
$y\in\{1,\dots,2000\}$, whether $4\,y!+1$ is a perfect square.  This is
decided exactly in integer arithmetic: maintaining $y!$ incrementally as an
exact (arbitrary-precision) integer, one computes
$s=\lfloor\sqrt{4\,y!+1}\rfloor$ by an exact integer square-root routine and
checks the integer identity $s^2=4\,y!+1$.  No floating-point arithmetic
enters, so the test is rigorous.  It returns ``square'' precisely for $y=2$
($4\cdot 2+1=9=3^2$, giving $x=1$) and $y=3$ ($4\cdot 6+1=25=5^2$, giving
$x=2$), and ``not a square'' for every other $y$ in the range.
\end{proof}

\section{Conclusion}

Collecting the rigorous statements:
\begin{itemize}
\item By Proposition~\ref{prop:square}, \eqref{eq:main} is equivalent to
``$4\,y!+1$ is a perfect square'', i.e.\ to \eqref{eq:m}.
\item By Proposition~\ref{prop:small}, the only solutions with $y\le 7$ are
$(1,2)$ and $(2,3)$, established by an explicit pronic-interval argument.
\item By Lemma~\ref{lem:coprime}, Lemma~\ref{lem:size}, and
Lemma~\ref{lem:bigprime}, any further solution has $x,x+1$ coprime, each
$>\sqrt{y!}-1$, with each prime power $p^{v_p(y!)}$ lying wholly in one
factor, and with every prime $p\in(y/2,y]$ exactly dividing one factor.
\item By Theorem~\ref{thm:finite}, \eqref{eq:main} has only finitely many
solutions; but this is ineffective (Remark~\ref{rem:ineff}).
\item By Proposition~\ref{prop:comp}, the only solutions with $y\le 2000$
are $(1,2)$ and $(2,3)$.
\item By Propositions~\ref{prop:nocong}--\ref{prop:nosize}, neither
congruence, nor valuation-parity, nor fixed-prime-power size arguments can
settle the problem.
\end{itemize}

\begin{remark}[The precise missing ingredient]\label{rem:missing}
A complete determination of the solution set would follow from any
\emph{effective} upper bound $y\le Y_0$ for solutions of \eqref{eq:main}
(equivalently of \eqref{eq:m}): one could then finish by the exact integer
test of Proposition~\ref{prop:comp} on $\{1,\dots,Y_0\}$, and for any
$Y_0\le 2000$ the determination is already complete. The Berend--Osgood
theorem supplies finiteness but no such $Y_0$, and no effective $Y_0$ is
known in the literature. Producing one is exactly the open part of the
problem. This is a genuine gap in the mathematical literature, not an
omission in the present exposition: the metatheorems of
Section~\ref{sec:obstruction} prove that the elementary tools (congruences,
parity, fixed-prime-power size) provably cannot supply such a bound.
\end{remark}

The honest conclusion is that the only \emph{known} solutions are
$(x,y)=(1,2)$ and $(x,y)=(2,3)$; that there is strong computational evidence
together with an (ineffective) finiteness theorem supporting the belief that
these are the only ones; but that a \emph{complete} determination is, with
present techniques, an \textbf{open problem} of Brocard--Ramanujan type. We
have proved rigorously everything that the current state of knowledge
permits, and we have pinpointed precisely which ingredient is missing,
namely an effective upper bound for $y$ in \eqref{eq:main}.
\end{document}
